% ===========================================================================
%  Chapitre 5 — Dérivation et sens de variation
%  PE 2026, composante ANALYSE
% ===========================================================================
\bmtchapitre{Dérivation et sens de variation}
  {Calculer la dérivée d'une fonction, y compris d'une fonction composée,
   étudier le signe de cette dérivée pour en déduire le sens de variation et
   dresser le tableau de variation, et écrire l'équation de la tangente à la
   courbe en un point donné.}
  {Analyse \textbullet\ environ 5 heures}

Le chapitre précédent disait ce qui se passe très loin ; celui-ci dit ce qui
se passe tout près. La dérivée mesure la pente de la courbe en un point : là
où elle est positive, la courbe monte ; là où elle est négative, la courbe
descend. Une inégalité sur $f'$ devient une flèche dans un tableau, et ce
tableau devient un dessin.

C'est le chapitre le plus rentable du livre. À l'examen, le calcul de $f'(x)$
et le tableau de variation qui en découle valent, à eux deux, entre trois et
quatre points sur les dix du problème — souvent davantage que le tracé.

% ---------------------------------------------------------------------------
\begin{revision}
Rien de nouveau : ce sont les acquis de première.

\begin{enumerate}[label=\textbf{\arabic*.}, leftmargin=*, itemsep=0.35em]
  \item Donner le coefficient directeur de la droite d'équation
        $y = -3x+7$.
  \item Écrire l'équation de la droite passant par $A(1\,;\,4)$ et de
        coefficient directeur $2$.
  \item Développer $(x-1)(x+3)$.
  \item Résoudre $-2x+6 > 0$.
  \item Donner le signe de $\dfrac{x-2}{(x+1)^{2}}$ selon les valeurs de
        $x$.
  \item Dans un tableau de variation, que signifie une flèche montante ?
\end{enumerate}
\end{revision}

\begin{automatismes}
Dérivées à sortir sans réfléchir.

\begin{enumerate}[label=\textbf{\alph*)}, leftmargin=*, itemsep=0.25em]
  \item Dérivée de $x^{2}$
  \item Dérivée de $x^{3}$
  \item Dérivée de $5x$
  \item Dérivée de $7$
  \item Dérivée de $\dfrac{1}{x}$
  \item Dérivée de $\sqrt{x}$
  \item Dérivée de $x^{2}-3x+1$
  \item Le signe de $(x+1)^{2}$
\end{enumerate}
\end{automatismes}

\begin{corrige}[automatismes]
a) $2x$ \quad b) $3x^{2}$ \quad c) $5$ \quad d) $0$ \quad
e) $-\frac{1}{x^{2}}$ \quad f) $\frac{1}{2\sqrt x}$ \quad g) $2x-3$ \quad
h) positif ou nul, nul seulement en $-1$.
\end{corrige}

% ===========================================================================
\section{Nombre dérivé et tangente}

\begin{definition}[nombre dérivé]
Soit $f$ définie sur un intervalle $I$ et $a \in I$. Si le quotient
\[
  \frac{f(a+h)-f(a)}{h}
\]
admet une limite finie quand $h$ tend vers $0$, cette limite s'appelle le
\textbf{nombre dérivé} de $f$ en $a$ et se note $f'(a)$.
\end{definition}

\begin{visualisation}[la tangente, c'est la pente]
\begin{center}
\begin{tikzpicture}
\begin{axis}[
  width = 10.4cm, height = 6.2cm,
  axis lines = middle, xlabel = {$x$}, ylabel = {$y$},
  xmin = -0.6, xmax = 4.4, ymin = -0.8, ymax = 5.2,
  xtick = {1,2,3,4}, ytick = {1,2,3,4,5},
  axis line style = {bmtGris},
  tick label style = {font=\scriptsize, color=bmtGris}, clip = true,
]
  \addplot[bmtIndigo, line width=1.4pt, domain=0:4, samples=140]
    {0.45*x^2-0.4*x+1};
  \addplot[bmtLaterite, line width=1.1pt, domain=0.6:3.4]
    {1.4*x-1.01};
  \filldraw[bmtVetiver] (axis cs:2,1.4) circle (2.1pt);
  \node[bmtVetiver, font=\scriptsize, anchor=north west]
    at (axis cs:2.05,1.3) {$A\left(a\,;\,f(a)\right)$};
  \node[bmtLaterite, font=\scriptsize, anchor=west]
    at (axis cs:3.05,3.5) {tangente : pente $f'(a)$};
  \draw[bmtGris, dash pattern=on 2pt off 2pt]
    (axis cs:2,0) -- (axis cs:2,1.4);
  \node[bmtGris, font=\scriptsize, anchor=north] at (axis cs:2,-0.05) {$a$};
\end{axis}
\end{tikzpicture}
\end{center}

Le nombre dérivé n'est pas une abstraction : c'est le coefficient directeur
de la droite qui « épouse » la courbe au point d'abscisse $a$. Une pente
positive, une courbe qui monte ; une pente nulle, une tangente horizontale,
et donc — le plus souvent — un sommet ou un creux.
\end{visualisation}

\begin{theoreme}[équation de la tangente]
Si $f$ est dérivable en $a$, la tangente à $\Cf$ au point d'abscisse $a$ a
pour équation
\[
  (T) : \quad y = f'(a)\,(x-a)+f(a).
\]
\end{theoreme}

\begin{methode}{écrire l'équation d'une tangente}
Trois nombres, dans cet ordre, et rien d'autre :
\begin{enumerate}[label=\textbf{\arabic*.}, leftmargin=*, itemsep=0.15em]
  \item calculer $f(a)$ ;
  \item calculer $f'(x)$, puis $f'(a)$ ;
  \item remplacer dans la formule et \emph{développer} pour obtenir la forme
        $y = mx+p$.
\end{enumerate}
\end{methode}

\begin{exemple}[une tangente complète]
Soit $f(x) = x^{2}-4x+5$ et $a = 3$.

$f(3) = 9-12+5 = 2$. Ensuite $f'(x) = 2x-4$, donc $f'(3) = 2$. La tangente a
pour équation
\[
  y = 2(x-3)+2 = 2x-4 .
\]
\end{exemple}

\begin{erreurclassique}
La formule est $y = f'(a)(x-a)+f(a)$, et non $y = f'(a)(x-a)+f'(a)$ ni
$y = f(a)(x-a)+f'(a)$. Moyen mnémotechnique : le nombre \emph{dérivé}
multiplie, le nombre \emph{image} s'ajoute.
\end{erreurclassique}

% ===========================================================================
\section{Les formules de dérivation}

\begin{propriete}[dérivées usuelles]
\begin{center}
\renewcommand{\arraystretch}{1.45}
\begin{tabular}{|L{3.9cm}|L{3.9cm}|L{3.5cm}|}
\hline
\textbf{Fonction} & \textbf{Dérivée} & \textbf{Sur} \\
\hline
$f(x) = k$ (constante) & $f'(x) = 0$ & $\R$ \\
$f(x) = x$ & $f'(x) = 1$ & $\R$ \\
$f(x) = x^{n}$, $n \geq 1$ & $f'(x) = n\,x^{n-1}$ & $\R$ \\
$f(x) = \dfrac1x$ & $f'(x) = -\dfrac{1}{x^{2}}$ & $\R^{*}$ \\
$f(x) = \sqrt x$ & $f'(x) = \dfrac{1}{2\sqrt x}$ & $\Rps$ \\
\hline
\end{tabular}
\end{center}
\end{propriete}

\begin{propriete}[opérations]
Si $u$ et $v$ sont dérivables :
\[
  (u+v)' = u'+v', \qquad (ku)' = k\,u', \qquad (uv)' = u'v+uv',
\]
\[
  \left(\frac1v\right)' = -\frac{v'}{v^{2}}, \qquad
  \left(\frac uv\right)' = \frac{u'v-uv'}{v^{2}}
  \quad (\text{là où } v \neq 0).
\]
\end{propriete}

\begin{propriete}[dérivée d'une fonction composée]
Si $g$ est dérivable en $u(x)$ et $u$ dérivable en $x$, alors
\[
  \left(g \circ u\right)'(x) = u'(x)\times g'\!\left(u(x)\right).
\]
Cas particuliers à connaître :
\[
  \left(u^{n}\right)' = n\,u'\,u^{n-1}, \qquad
  \left(\sqrt u\right)' = \frac{u'}{2\sqrt u}.
\]
\end{propriete}

\begin{exemple}[quatre dérivées, quatre règles]
\textbf{a)} $f(x) = (3x-1)^{4}$ : ici $u = 3x-1$ et $u' = 3$, donc
$f'(x) = 4\times3\times(3x-1)^{3} = 12(3x-1)^{3}$.

\textbf{b)} $g(x) = \dfrac{2x+1}{x-3}$ : avec $u = 2x+1$, $v = x-3$,
\[
  g'(x) = \frac{2(x-3)-(2x+1)\times1}{(x-3)^{2}}
  = \frac{-7}{(x-3)^{2}} .
\]

\textbf{c)} $h(x) = x^{2}\sqrt x$ — on écrit plutôt
$h(x) = x^{2}\times x^{1/2}$, ou l'on applique la règle du produit :
$h'(x) = 2x\sqrt x+x^{2}\times\frac{1}{2\sqrt x}
= 2x\sqrt x+\frac{x\sqrt x}{2} = \frac{5x\sqrt x}{2}$.

\textbf{d)} $k(x) = \sqrt{x^{2}+1}$ : $u = x^{2}+1$, $u' = 2x$, donc
$k'(x) = \dfrac{2x}{2\sqrt{x^{2}+1}} = \dfrac{x}{\sqrt{x^{2}+1}}$.
\end{exemple}

\begin{avertissement}
Après un calcul de dérivée par la règle du quotient, on \textbf{développe le
numérateur seulement} et l'on \textbf{laisse le dénominateur factorisé} : il
est carré, donc positif, et c'est précisément ce qui permettra de lire le
signe de $f'$ sur le seul numérateur.
\end{avertissement}

\begin{entrainement}[calculs de dérivées]
\exo[\niveau{1}] $f(x) = 3x^{2}-5x+2$

\exo[\niveau{1}] $f(x) = x^{3}-\dfrac{1}{x}$

\exo[\niveau{2}] $f(x) = (2x+5)^{3}$

\exo[\niveau{2}] $f(x) = \dfrac{x+1}{x-1}$

\exo[\niveau{2}] $f(x) = \dfrac{3x}{3+4x}$

\exo[\niveau{2}] $f(x) = x+\dfrac{1}{x+1}$

\exo[\niveau{3}] $f(x) = \dfrac{x^{2}-3x+4}{x-2}$

\exo[\niveau{3}] $f(x) = \sqrt{4x-1}$
\end{entrainement}

\begin{corrige}[calculs de dérivées]
\textbf{1.} $6x-5$.
\textbf{2.} $3x^{2}+\frac{1}{x^{2}}$.
\textbf{3.} $6(2x+5)^{2}$.
\textbf{4.} $\frac{(x-1)-(x+1)}{(x-1)^{2}} = \frac{-2}{(x-1)^{2}}$.
\textbf{5.} $\frac{3(3+4x)-3x\times4}{(3+4x)^{2}} = \frac{9}{(3+4x)^{2}}$.
\textbf{6.} $1-\frac{1}{(x+1)^{2}} = \frac{(x+1)^{2}-1}{(x+1)^{2}}
= \frac{x(x+2)}{(x+1)^{2}}$.
\textbf{7.} $\frac{(2x-3)(x-2)-(x^{2}-3x+4)}{(x-2)^{2}}
= \frac{x^{2}-4x+2}{(x-2)^{2}}$.
\textbf{8.} $\frac{4}{2\sqrt{4x-1}} = \frac{2}{\sqrt{4x-1}}$.
\end{corrige}

% ===========================================================================
\section{Du signe de la dérivée au tableau de variation}

\begin{theoreme}[variations et signe de la dérivée]
Soit $f$ dérivable sur un intervalle $I$.
\begin{itemize}[leftmargin=1.3em, itemsep=0.15em]
  \item Si $f'(x) > 0$ sur $I$, alors $f$ est strictement croissante sur
        $I$.
  \item Si $f'(x) < 0$ sur $I$, alors $f$ est strictement décroissante sur
        $I$.
  \item Si $f'(x) = 0$ sur $I$, alors $f$ est constante sur $I$.
\end{itemize}
\end{theoreme}

\begin{propriete}[extremum]
Si $f'$ s'annule en $a$ \textbf{en changeant de signe}, alors $f$ admet un
extremum local en $a$ : un maximum si $f'$ passe du positif au négatif, un
minimum si elle passe du négatif au positif.
\end{propriete}

\begin{avertissement}
« $f'(a) = 0$ » ne suffit pas à garantir un extremum : il faut le
\emph{changement de signe}. La fonction $f(x) = x^{3}$ a une dérivée nulle en
$0$, et pourtant elle est croissante sur $\R$ tout entier — sa courbe
présente en $0$ une tangente horizontale sans sommet ni creux.
\end{avertissement}

\begin{methode}{dresser un tableau de variation}
\begin{enumerate}[label=\textbf{\arabic*.}, leftmargin=*, itemsep=0.15em]
  \item Écrire l'ensemble de définition et le rappeler en première ligne du
        tableau, avec les valeurs interdites comme colonnes.
  \item Calculer $f'(x)$ et le mettre sous la forme
        $\dfrac{\text{numérateur}}{\text{carré}}$ quand c'est un quotient.
  \item Chercher les racines du numérateur et faire une ligne de signes.
  \item Reporter les flèches, puis remplir les extrémités avec les limites
        calculées au chapitre précédent — c'est ce qui rapporte les points.
\end{enumerate}
\end{methode}

\begin{exemple}[une étude conduite jusqu'au tableau]
Soit $f(x) = x+\dfrac{1}{x+1}$ sur $\R\setminus\{-1\}$.

On a calculé $f'(x) = \dfrac{x(x+2)}{(x+1)^{2}}$. Le dénominateur est
strictement positif ; le signe de $f'$ est donc celui de $x(x+2)$, trinôme de
racines $-2$ et $0$, à coefficient dominant positif : positif à l'extérieur
de $\intervalle{-2}{0}$, négatif entre.

\begin{center}\footnotesize
\renewcommand{\arraystretch}{1.3}
\begin{tabular}{|c|ccccccccc|}
\hline
$x$ & $-\infty$ & & $-2$ & & $-1$ & & $0$ & & $+\infty$\\
\hline
$f'(x)$ & & $+$ & $0$ & $-$ & \textbardbl & $-$ & $0$ & $+$ & \\
\hline
$f$ & $-\infty$ & $\nearrow$ & $-3$ & $\searrow$ & \textbardbl &
$\searrow$ & $1$ & $\nearrow$ & $+\infty$\\
\hline
\end{tabular}
\end{center}

La double barre marque la valeur interdite $-1$ : la fonction n'y est pas
définie, et le tableau se coupe en deux. Le maximum local vaut
$f(-2) = -2-1 = -3$, le minimum local $f(0) = 0+1 = 1$ — un maximum plus bas
que le minimum, ce qui n'a rien de contradictoire puisqu'ils appartiennent à
deux branches séparées.
\end{exemple}

\begin{entrainement}[variations]
\exo[\niveau{1}] $f(x) = x^{2}-6x+5$ : dresser le tableau de variation sur
$\R$.

\exo[\niveau{2}] $f(x) = -x^{3}+3x$ : étudier les variations sur $\R$ et
donner les extremums.

\exo[\niveau{2}] $f(x) = \dfrac{2x+1}{x-3}$ : montrer que $f$ est
strictement décroissante sur chacun des deux intervalles de son ensemble de
définition.

\exo[\niveau{3}] $f(x) = \dfrac{x^{2}+1}{x}$ : dresser le tableau de
variation complet sur $\R^{*}$, limites comprises.

\exo[\niveau{3}] $f(x) = \dfrac{x^{2}-3x+4}{x-2}$ : montrer que $f'$ s'annule
en $2-\sqrt2$ et $2+\sqrt2$, et en déduire le tableau de variation.
\end{entrainement}

\begin{corrige}[variations]
\textbf{1.} $f'(x) = 2x-6$, négative avant $3$, positive après : $f$ décroît
sur $\intervalleof{-\infty}{3}$, croît sur $\intervallefo{3}{+\infty}$,
minimum $f(3) = -4$.
\textbf{2.} $f'(x) = -3x^{2}+3 = -3(x-1)(x+1)$ : négative à l'extérieur de
$\intervalle{-1}{1}$, positive à l'intérieur. Minimum local $f(-1) = -2$,
maximum local $f(1) = 2$.
\textbf{3.} $f'(x) = \frac{2(x-3)-(2x+1)}{(x-3)^{2}} = \frac{-7}{(x-3)^{2}}
< 0$ : la fonction décroît sur $\intervalleo{-\infty}{3}$ et sur
$\intervalleo{3}{+\infty}$ — mais pas sur la réunion des deux.
\textbf{4.} $f(x) = x+\frac1x$ et $f'(x) = 1-\frac{1}{x^{2}}
= \frac{x^{2}-1}{x^{2}}$ : positive à l'extérieur de $\intervalle{-1}{1}$,
négative à l'intérieur. Maximum local $f(-1) = -2$, minimum local
$f(1) = 2$ ; limites $-\infty$ et $+\infty$ aux bornes infinies, $-\infty$
en $0^{-}$ et $+\infty$ en $0^{+}$.
\textbf{5.} $f'(x) = \frac{x^{2}-4x+2}{(x-2)^{2}}$ ; le numérateur a pour
discriminant $16-8 = 8$ et pour racines $2\pm\sqrt2$. La dérivée est donc
positive à l'extérieur de $\intervalle{2-\sqrt2}{2+\sqrt2}$ et négative à
l'intérieur (valeur interdite $2$ exclue).
\end{corrige}

\begin{qcm}
\exo La dérivée de $\dfrac{1}{x^{3}}$ est :
\textbf{a)} $\dfrac{1}{3x^{2}}$ \quad \textbf{b)} $-\dfrac{3}{x^{4}}$ \quad
\textbf{c)} $-\dfrac{1}{3x^{2}}$ \quad \textbf{d)} $3x^{2}$
\reponseqcm{b}

\exo Si $f'(x) < 0$ sur $\intervalle{0}{5}$, alors sur cet intervalle $f$
est :
\textbf{a)} croissante \quad \textbf{b)} décroissante \quad
\textbf{c)} constante \quad \textbf{d)} négative \reponseqcm{b}

\exo La tangente à $\Cf$ au point d'abscisse $2$ est horizontale si :
\textbf{a)} $f(2) = 0$ \quad \textbf{b)} $f'(2) = 0$ \quad
\textbf{c)} $f'(2) = 2$ \quad \textbf{d)} $f(2) = 2$ \reponseqcm{b}

\exo La dérivée de $(x^{2}+1)^{5}$ est :
\textbf{a)} $5(x^{2}+1)^{4}$ \quad \textbf{b)} $10x(x^{2}+1)^{4}$ \quad
\textbf{c)} $2x(x^{2}+1)^{5}$ \quad \textbf{d)} $10(x^{2}+1)^{4}$
\reponseqcm{b}

\exo Si $f'$ s'annule en $a$ sans changer de signe, alors en $a$ la courbe
admet :
\textbf{a)} un maximum \quad \textbf{b)} un minimum \quad
\textbf{c)} une tangente horizontale sans extremum \quad
\textbf{d)} une asymptote \reponseqcm{c}
\end{qcm}

% ===========================================================================
\begin{annales}
Comme au chapitre précédent, le problème de la série~A construit ses dérivées
sur le logarithme ou l'exponentielle ; les énoncés ci-dessous sont donc des
exercices réels portant sur des fonctions rationnelles ou polynômes, tirés du
même fonds vérifié.

\exoann{Baccalauréat probatoire, Cameroun, série D, session 2022 —
exercice 4 (extrait)}
Soit $f(x) = \dfrac{3x}{3+4x}$.
\begin{enumerate}[label=\textbf{\alph*)}, leftmargin=*, itemsep=0.15em]
  \item Montrer que $f'(x) = \dfrac{9}{(3+4x)^{2}}$.
  \item En déduire le sens de variation de $f$ sur chacun des intervalles
        de son ensemble de définition.
\end{enumerate}

\exoann{Baccalauréat probatoire, Cameroun, séries D et TI, session 2023 —
exercice 4 (extrait)}
Soit $f(x) = x+\dfrac{1}{x+1}$ sur $\R\setminus\{-1\}$.
\begin{enumerate}[label=\textbf{\alph*)}, leftmargin=*, itemsep=0.15em]
  \item Montrer que $f'(x) = \dfrac{x(x+2)}{(x+1)^{2}}$.
  \item Étudier le signe de $f'(x)$ et dresser le tableau de variation.
\end{enumerate}

\exoann{Baccalauréat probatoire, Cameroun, série IH, session 2023 —
problème (extrait)}
La courbe de la fonction $f(x) = ax^{2}+bx-3$ passe par les points
$A(-2\,;\,-5)$ et $B(-1\,;\,-6)$.
\begin{enumerate}[label=\textbf{\alph*)}, leftmargin=*, itemsep=0.15em]
  \item Déterminer les réels $a$ et $b$.
  \item Écrire une équation de la tangente à la courbe au point $A$.
\end{enumerate}

\exoann{Baccalauréat blanc, Lycée Philibert Tsiranana, session 2022-2023 —
problème (extrait)}
Soit $f(x) = \dfrac{x^{2}-3x+4}{x-2}$. Montrer que $f$ admet deux extremums,
atteints en $2-\sqrt2$ et $2+\sqrt2$, et préciser leur nature.
\end{annales}

\begin{corrige}[annales]
\textbf{1. a)} $f'(x) = \frac{3(3+4x)-3x\times4}{(3+4x)^{2}}
= \frac{9+12x-12x}{(3+4x)^{2}} = \frac{9}{(3+4x)^{2}}$.
\textbf{b)} Ce quotient est strictement positif partout où il est défini :
$f$ est strictement croissante sur $\intervalleo{-\infty}{-\frac34}$ et sur
$\intervalleo{-\frac34}{+\infty}$.

\textbf{2. a)} $f'(x) = 1-\frac{1}{(x+1)^{2}} = \frac{(x+1)^{2}-1}{(x+1)^{2}}
= \frac{x^{2}+2x}{(x+1)^{2}} = \frac{x(x+2)}{(x+1)^{2}}$.
\textbf{b)} Le dénominateur est positif ; $x(x+2)$ est positif à l'extérieur
de $\intervalle{-2}{0}$. Donc $f$ croît sur $\intervalleof{-\infty}{-2}$,
décroît sur $\intervallefo{-2}{-1}$ puis sur $\intervalleof{-1}{0}$, et croît
sur $\intervallefo{0}{+\infty}$. Maximum local $f(-2) = -3$, minimum local
$f(0) = 1$.

\textbf{3. a)} $A$ donne $4a-2b-3 = -5$, soit $4a-2b = -2$ ; $B$ donne
$a-b-3 = -6$, soit $a-b = -3$. En résolvant : $a = 2$ et $b = 5$.
\textbf{b)} $f(x) = 2x^{2}+5x-3$, $f'(x) = 4x+5$, donc $f'(-2) = -3$, et
$f(-2) = -5$. La tangente a pour équation $y = -3(x+2)-5 = -3x-11$.

\textbf{4.} $f'(x) = \frac{x^{2}-4x+2}{(x-2)^{2}}$, dont le numérateur a pour
racines $2\pm\sqrt2$. Le signe de $f'$ est celui de ce trinôme : positif,
puis négatif, puis positif. En $2-\sqrt2 \approx 0{,}59$, $f$ passe de
croissante à décroissante : c'est un maximum local. En
$2+\sqrt2 \approx 3{,}41$, c'est un minimum local.
\end{corrige}

% ===========================================================================
\begin{vraifaux}
\exo Si $f$ est croissante sur $I$, alors $f'(x) > 0$ pour tout $x$ de $I$.

\exo Une fonction dont la dérivée est constante est une fonction affine.

\exo La dérivée d'un produit est le produit des dérivées.

\exo Si $f'(2) = 0$, la courbe admet un extremum en $2$.

\exo La tangente à une courbe en un point peut couper cette courbe.
\end{vraifaux}

\begin{corrige}[vrai ou faux]
\textbf{1.} Faux — elle est \emph{positive ou nulle} ; $x \mapsto x^{3}$ est
croissante et sa dérivée s'annule en $0$.
\textbf{2.} Vrai — si $f' = m$, alors $f(x) = mx+p$.
\textbf{3.} Faux — $(uv)' = u'v+uv'$.
\textbf{4.} Faux — il faut que $f'$ change de signe.
\textbf{5.} Vrai — la tangente épouse la courbe localement, mais rien ne
l'empêche de la traverser, notamment en un point d'inflexion.
\end{corrige}

% ===========================================================================
\begin{situation}[le rendement d'une rizière]
Un agronome modélise le rendement $R$ d'une rizière, en tonnes par hectare,
en fonction de la quantité $x$ d'engrais épandue, en quintaux par hectare,
par
\[
  R(x) = -0{,}2x^{2}+2x+3, \qquad 0 \leq x \leq 8 .
\]

\begin{enumerate}[label=\textbf{\arabic*.}, leftmargin=*, itemsep=0.3em]
  \item Calculer le rendement sans engrais.
  \item Calculer $R'(x)$, étudier son signe sur $\intervalle{0}{8}$ et
        dresser le tableau de variation de $R$.
  \item Quelle quantité d'engrais donne le rendement maximal ? Que vaut ce
        rendement ?
  \item Que se passe-t-il si l'on épand $8$ quintaux ? Commenter en une
        phrase ce que dit le modèle.
\end{enumerate}
\end{situation}

\begin{corrige}[le rendement d'une rizière]
\textbf{1.} $R(0) = 3$ tonnes par hectare.

\textbf{2.} $R'(x) = -0{,}4x+2$, positive pour $x < 5$, nulle en $5$,
négative pour $x > 5$. $R$ croît sur $\intervalle{0}{5}$, décroît sur
$\intervalle{5}{8}$.

\textbf{3.} Le maximum est atteint pour $x = 5$ quintaux, et vaut
$R(5) = -5+10+3 = 8$ tonnes par hectare.

\textbf{4.} $R(8) = -12{,}8+16+3 = 6{,}2$ tonnes : au-delà de cinq quintaux,
l'engrais supplémentaire fait \emph{baisser} le rendement. Le modèle traduit
un excès d'azote — et rappelle qu'un maximum n'est pas au bout de
l'intervalle, mais là où la dérivée s'annule.
\end{corrige}

% ===========================================================================
\begin{jemeteste}
Vingt-cinq minutes.

\begin{enumerate}[label=\textbf{\arabic*.}, leftmargin=*, itemsep=0.3em]
  \item Dériver $f(x) = 4x^{3}-2x^{2}+7$.
  \item Dériver $g(x) = \dfrac{x-2}{x+2}$ et donner le signe de $g'$.
  \item Dériver $h(x) = (1-3x)^{4}$.
  \item Écrire l'équation de la tangente à la courbe de
        $f(x) = x^{2}-x$ au point d'abscisse $-1$.
  \item Dresser le tableau de variation de $k(x) = x^{3}-3x^{2}$ sur $\R$.
\end{enumerate}
\end{jemeteste}

\begin{corrige}[je me teste]
\textbf{1.} $12x^{2}-4x$.
\textbf{2.} $g'(x) = \frac{(x+2)-(x-2)}{(x+2)^{2}} = \frac{4}{(x+2)^{2}}
> 0$ : $g$ croît sur chacun des deux intervalles.
\textbf{3.} $-12(1-3x)^{3}$.
\textbf{4.} $f(-1) = 2$, $f'(x) = 2x-1$ donc $f'(-1) = -3$ ; la tangente a
pour équation $y = -3(x+1)+2 = -3x-1$.
\textbf{5.} $k'(x) = 3x^{2}-6x = 3x(x-2)$ : $k$ croît sur
$\intervalleof{-\infty}{0}$, décroît sur $\intervalle{0}{2}$, croît sur
$\intervallefo{2}{+\infty}$. Maximum local $k(0) = 0$, minimum local
$k(2) = -4$.
\end{corrige}

% ===========================================================================
\begin{commeaubac}
\bmtsource{Exercice original au format de l'épreuve — série A \textbullet\ 5 points}
Soit $f$ la fonction définie sur $\R\setminus\{2\}$ par
$f(x) = \dfrac{x^{2}}{x-2}$, de courbe $\Cf$.

\begin{enumerate}[label=\textbf{\arabic*.}, leftmargin=*, itemsep=0.25em]
  \item Montrer que, pour tout $x \neq 2$,
        $f'(x) = \dfrac{x(x-4)}{(x-2)^{2}}$. \bareme{1,5 pt}
  \item Étudier le signe de $f'(x)$ sur $\R\setminus\{2\}$. \bareme{1 pt}
  \item Calculer $f(0)$ et $f(4)$, puis dresser le tableau de variation de
        $f$ (les limites ne sont pas demandées). \bareme{1,5 pt}
  \item Écrire une équation de la tangente $(T)$ à $\Cf$ au point d'abscisse
        $1$. \bareme{1 pt}
\end{enumerate}
\end{commeaubac}

\begin{corrige}[comme au baccalauréat]
\textbf{1.} Avec $u = x^{2}$ et $v = x-2$ :
\[
  f'(x) = \frac{2x(x-2)-x^{2}}{(x-2)^{2}}
  = \frac{2x^{2}-4x-x^{2}}{(x-2)^{2}} = \frac{x^{2}-4x}{(x-2)^{2}}
  = \frac{x(x-4)}{(x-2)^{2}} .
\]

\textbf{2.} Le dénominateur est strictement positif. Le numérateur $x(x-4)$
est positif à l'extérieur de $\intervalle{0}{4}$, négatif à l'intérieur.
Donc $f' > 0$ sur $\intervalleo{-\infty}{0}$ et sur
$\intervalleo{4}{+\infty}$, et $f' < 0$ sur $\intervalleo{0}{2}$ et
$\intervalleo{2}{4}$.

\textbf{3.} $f(0) = 0$ et $f(4) = \frac{16}{2} = 8$. La fonction croît
jusqu'en $0$ (maximum local $0$), décroît jusqu'à la valeur interdite $2$,
décroît encore de $2$ à $4$ (minimum local $8$), puis croît. La double barre
en $2$ sépare le tableau en deux.

\textbf{4.} $f(1) = \frac{1}{-1} = -1$ et
$f'(1) = \frac{1\times(-3)}{1} = -3$. D'où
$(T) : y = -3(x-1)-1 = -3x+2$.
\end{corrige}

\begin{notepourlaclasse}
Cinq heures, dont au moins deux sur le seul passage « signe de $f'$
$\rightarrow$ tableau ». Les élèves savent dériver bien avant de savoir
conclure : c'est la conclusion qui rapporte les points.

Insistez sur l'écriture $f'(x) = \frac{N(x)}{(\cdots)^{2}}$ et sur la phrase
qui l'accompagne : « le dénominateur étant un carré, il est strictement
positif ; le signe de $f'$ est donc celui de son numérateur ». Cette phrase
figure telle quelle dans les corrigés officiels et vaut souvent un demi-point
à elle seule.

L'avertissement sur $x \mapsto x^{3}$ mérite d'être fait au tableau plutôt
que lu : tracez la courbe, montrez la tangente horizontale en $0$, et
demandez où est le sommet. C'est la seule façon de faire tomber l'idée que
« dérivée nulle » signifie « extremum ».
\end{notepourlaclasse}
